% Capitulo: A estrutura dos times de dados
\chapter{A estrutura dos times de dados}
\label{cap:estrutura-times}

\rascunho{
  \textbf{Ideia do capítulo.} Abrir a Parte~\ref{parte:agil} mostrando quem compõe um
  time de dados. A estrutura tradicional tem um time de engenheiros com seu gerente, um
  time de analistas com seu gerente, às vezes alguém de ágil, um product manager, e isso
  sem contar os times de negócio. Fazer um paralelo de como essa estrutura pode ser
  reduzida com a IA e qual a minha previsão para o futuro.
  \begin{itemize}
    \item A Figura~\ref{fig:estrutura-times} mostra as duas estruturas lado a lado.
    \item Lei de Conway: a arquitetura dos sistemas tende a copiar a estrutura de
      comunicação de quem os constrói \parencite{conway1968}. Em dados, dois times
      separados costumam gerar duas camadas que não conversam.
    \item Referências para formas de organizar times: \textcite{skelton2019} e, para
      times orientados a domínio, \textcite{dehghani2022}.
  \end{itemize}
}

% ---------------------------------------------------------------------------
\section{A estrutura tradicional}
\label{sec:estrutura-tradicional}

\rascunho{
  \begin{itemize}
    \item Cada papel e o que faz: head de dados, gerente de engenharia e engenheiros,
      gerente de analytics e analistas, product manager, agilista (scrum master ou agile
      coach), e do outro lado os times de negócio com seus analistas e gestores.
    \item Por que ela surgiu: as ferramentas eram especializadas e difíceis, e cada etapa
      (extrair, modelar, visualizar) pedia gente dedicada.
    \item O custo: muitas passagens de bastão. Um pedido de negócio passa pelo product
      manager, vira história, vai para a engenharia, volta para os analistas, e cada
      passagem perde contexto e tempo.
    \item Relato: as estruturas que eu vi ou montei. \pendente{escrever com relato, sem nomes}
  \end{itemize}
}

\pendente{escrever}

\begin{figure}[htbp]
  \centering
  % Estrutura tradicional de um time de dados e uma versao enxuta possivel (cap. 23).
\begin{tikzpicture}[diagrama, x=1mm, y=1mm,
    titulo/.style={anchor=west, font=\footnotesize\scshape, text=cinza},
    cargo/.style={caixa, minimum width=24mm, minimum height=8mm, font=\scriptsize},
    pilha/.style={cargo, double copy shadow={shadow xshift=1mm, shadow yshift=-1mm, draw=tinta, fill=white}},
    negocio/.style={caixa, draw=cinza, dashed, font=\scriptsize, text=cinza},
    hierarquia/.style={draw=tinta, line width=0.5pt}]

  % --- estrutura tradicional
  \node[titulo] at (0,58) {estrutura tradicional};
  \node[cargo] (head) at (58,50) {Head de dados};
  \node[cargo] (ge) at (14,37) {Gerente de\\engenharia};
  \node[cargo] (ga) at (44,37) {Gerente de\\analytics};
  \node[cargo] (pm) at (74,37) {Product\\manager};
  \node[cargo] (ag) at (103,37) {Agilista\\(scrum master)};
  \foreach \n in {ge,ga,pm,ag} \draw[hierarquia] (head.south) -- ++(0,-3) -| (\n.north);
  \node[pilha] (eng) at (14,23) {Engenheiros\\de dados};
  \node[pilha] (ana) at (44,23) {Analistas de\\dados e BI};
  \draw[hierarquia] (ge) -- (eng);
  \draw[hierarquia] (ga) -- (ana);
  \node[negocio, minimum width=50mm, minimum height=8mm] (neg1) at (88,23)
    {Times de negócio: analistas de\\negócio, gestores, usuários};
  \draw[ligacao, dashed] (pm.south) -- (pm.south |- neg1.north);

  % --- versao enxuta
  \draw[fio] (0,12) -- (118,12);
  \node[titulo] at (0,7) {uma versão enxuta possível};
  \node[destaquecaixa, minimum width=34mm, minimum height=8mm, font=\scriptsize] (lider) at (58,0)
    {Líder de dados\\(produto e fluxo)};
  \node[pilha] (prof) at (40,-14) {Profissionais de\\dados (ponta a ponta)};
  \node[caixa, draw=cinza, minimum width=24mm, minimum height=8mm, font=\scriptsize] (ia) at (76,-14)
    {Agentes de IA\\(revisados pelo time)};
  \draw[hierarquia] (lider.south) -- ++(0,-3) -| (prof.north);
  \draw[ligacao] (lider.south) -- ++(0,-3) -| (ia.north);
  \node[negocio, minimum width=26mm, minimum height=8mm] (neg2) at (105,-14)
    {Negócio, com\\analistas próprios};
  \draw[ligacao, dashed] (ia.east) -- (neg2.west);
\end{tikzpicture}

  \caption{A estrutura tradicional de um time de dados e uma versão enxuta possível com
    a IA. Em linha tracejada, a relação com o negócio.}
  \label{fig:estrutura-times}
  \fonte{elaborado pelo autor.}
\end{figure}

% ---------------------------------------------------------------------------
\section{O que a IA comprime}
\label{sec:ia-comprime}

\rascunho{
  \begin{itemize}
    \item A fronteira entre engenharia e análise ficou mais fina: o mesmo profissional
      consegue ir da ingestão ao dashboard com a IA ajudando nas partes que não domina
      (ligação com o analytics engineering do capítulo~\ref{cap:nuvem}).
    \item Papéis que se juntam: com times menores, o líder de dados acumula a gestão de
      produto e do fluxo; a função de agilista vira prática do time, não um cargo.
    \item O que não comprime: conhecer o negócio (capítulo~\ref{cap:negocio}), revisar o
      que a IA produz (capítulo~\ref{cap:programar}), decidir o que construir e com que
      frequência (capítulo~\ref{cap:frequencia}).
    \item O risco: reduzir a estrutura sem reduzir a demanda, e transformar o time enxuto
      num time sobrecarregado. Ou cortar os juniores e ficar sem quem vire sênior daqui a
      alguns anos.
  \end{itemize}
}

\pendente{escrever}

% ---------------------------------------------------------------------------
\section{Uma estrutura enxuta}
\label{sec:estrutura-enxuta}

\rascunho{
  \begin{itemize}
    \item Uma possibilidade (parte de baixo da Figura~\ref{fig:estrutura-times}): um líder
      de dados que cuida de produto e fluxo, profissionais de dados que trabalham de ponta
      a ponta, agentes de IA executando tarefas revisadas pelo time, e o negócio com
      analistas próprios, mais perto dos dados.
    \item Como o Scrum com Kanban dos próximos capítulos se adapta a um time de três ou
      quatro pessoas.
    \item O que fica com uma plataforma central e o que vai para os domínios de negócio
      \parencite{dehghani2022,skelton2019}.
  \end{itemize}
}

\pendente{escrever}

% ---------------------------------------------------------------------------
\section{Uma previsão para o futuro}
\label{sec:previsao-times}

\rascunho{
  A previsão é minha e deve ser escrita como aposta, com as razões. Ideias para
  escolher ou descartar:
  \begin{itemize}
    \item Times de dados menores e mais seniores, com a IA fazendo parte do trabalho que
      era de pessoas em início de carreira.
    \item O analista de negócio e o analista de dados convergindo num mesmo papel, mais
      perto da área de negócio.
    \item O gerente de dados como gestor de produto, de fluxo e de consumo (créditos e
      tokens, capítulo~\ref{cap:custo-consumo}) ao mesmo tempo.
    \item A pergunta em aberto: como se forma um sênior se o trabalho de júnior encolher?
  \end{itemize}
  \pendente{escrever a previsão; alinhar com o epílogo para não repetir}
}

\pendente{escrever}

% ---------------------------------------------------------------------------
\begin{leituras}
  \leitura{skelton2019}{Um vocabulário claro para tipos de time (de fluxo, de plataforma,
    de apoio) e para as formas de interação entre eles.}
  \leitura{dehghani2022}{A proposta de organizar dados por domínio de negócio, com times
    responsáveis pelos próprios produtos de dados.}
  \leitura{conway1968}{O artigo curto que deu origem à lei de Conway, ainda útil para
    entender por que a estrutura do time aparece na arquitetura.}
\end{leituras}
